%abstract traces
\newcommand{\noEvent}[1]{\stackrel{#1}{\ensuremath{\cdot\cdot}}}
\newcommand{\anyEvent}[0]{\ensuremath\noEvent{}\xspace}
%chop
\newcommand{\chop}{\ensuremath{*\!*}\xspace}
%concat
\newcommand{\concat}{\ensuremath\cdot}
%mu formula
\newcommand{\muFml}[2]{\ensuremath{\mu#1.#2}}
%state formula
\newcommand{\stateFml}[1]{\lceil#1\rceil}
%events
\newcommand{\pushEv}{\ensuremath{\mathsf{push}}\xspace}
\newcommand{\popEv}{\ensuremath{\mathsf{pop}}\xspace}
\newcommand{\callEv}{\ensuremath{\mathsf{call}}\xspace}
\newcommand{\retEv}{\ensuremath{\mathsf{ret}}\xspace}
\newcommand{\startEv}{\ensuremath{\mathsf{start}}\xspace}
\newcommand{\finishEv}{\ensuremath{\mathsf{finish}}\xspace}
\newcommand{\invocEv}{\ensuremath{\mathsf{invoc}}\xspace}
\newcommand{\invocREv}{\ensuremath{\mathsf{invocR}}\xspace}
\newcommand{\semchop}{\ensuremath{*\!*}\xspace}

% TODO figure out what \xasync is
\newcommand{\xasync}[1]{#1}

%observations
\newcommand{\obs}[2]{\mho\, \xasync{#1}~\mathbf{as}~#2}
\newcommand{\as}[2]{\xasync{#1}~\mathbf{as}~#2}
\newcommand{\asMany}[2]{\overline{\xasync{#1}}~\mathbf{as}~\overline{#2}}
\newcommand{\obsP}[3]{\mho\, \xasync{#1}~\mathbf{as}~#2.\stateFml{#3}}
\newcommand{\obsPMany}[2]{\mho\, \overline{\xasync{#1}}~\mathbf{as}~\overline{#2}.}
\newcommand{\obsPManyL}[2]{\mho\, \overline{#1}~\mathbf{as}~\overline{#2}.}

%cats
\newcommand{\assume}[1]{\ensuremath{<\hspace{-5pt}<\!#1}}
\newcommand{\continue}[1]{\ensuremath{#1\!>\hspace{-5pt}>}}
\newcommand{\CAT}[3]{\ensuremath{\assume{#1} \ | \ #2 \ | \ \continue{#3}}}
\newcommand{\CATMulti}[3]{\begin{align*}\ensuremath{&\assume{#1} \ | \\  &#2 \ | \\ &\continue{#3}\end{align*}}}
\newcommand{\CATShort}[3]{\ensuremath{\assume{#1} \big|\big. #2 \big|\big. \continue{#3}}}

%misc
\newcommand{\idm}[1]{\ensuremath \mathsf{#1}\xspace}
\newcommand{\fixpoint}[2]{\ensuremath(\mu#1.#2)}
\newcommand{\nomfsm}[0]{\noEvent{\mfsm}}
\newcommand{\nomfsmExcept}[1]{\noEvent{\mfsm\setminus\{#1\}}}
